% Figura: interpretación geométrica del gradiente
\begin{figure}[ht]
\centering
\begin{minipage}[t]{0.47\textwidth}\centering
\begin{tikzpicture}[>=Stealth, baseline=(current bounding box.north)]
  \coordinate (P) at (0,0);
  \draw[vect, line width=1.4pt, acento] (P) -- (4.0,1.6) node[above right, text=acento] {$\nab f$};
  \draw[vect, line width=1.2pt, azulprof] (P) -- (2.2,-0.9) node[right, text=azulprof] {$\dl$};
  % ángulo
  \draw[grispizarra, line width=0.5pt] (1.1,{1.1*1.6/4.0}) arc (21.8:-22.2:1.1);
  \node[font=\small] at (1.45,0.0) {$\theta$};
  % proyección de dl sobre grad f
  \pgfmathsetmacro{\pr}{(2.2*4.0-0.9*1.6)/(4.0*4.0+1.6*1.6)}
  \draw[densely dashed, grispizarra] (2.2,-0.9) -- ({\pr*4.0},{\pr*1.6});
  \draw[line width=2.4pt, azulclaro] (P) -- ({\pr*4.0},{\pr*1.6});
  \node[font=\footnotesize, text=azulacero, rotate=21.8, anchor=south] at ({\pr*2.0},{\pr*0.8}) {$\dl\cos\theta$};
  \draw[grispizarra, line width=0.5pt] ({\pr*4.0-0.18*0.928},{\pr*1.6-0.18*0.371})
     -- ++(0.371*0.18,-0.928*0.18) -- ++(0.928*0.18,0.371*0.18);
  \fill (P) circle (1.6pt);
\end{tikzpicture}\\[4pt]{\small\scshape (a)}\; $\dif f=\norm{\nab f}\,\norm{\dl}\cos\theta$
\end{minipage}\hfill
\begin{minipage}[t]{0.47\textwidth}\centering
\begin{tikzpicture}[>=Stealth, baseline=(current bounding box.north), scale=0.9, transform shape]
  \foreach \c/\pc in {1/100,2/65,3/35}{%
    \draw[azulprof!\pc!azulclaro, line width=0.9pt] (0,0) ellipse ({1.0*\c} and {0.5*\c});}
  \fill[azulprof] (0,0) circle (1.4pt);
  
  % gradientes en puntos de las elipses (apuntan hacia el centro)
  \foreach \c/\t in {2/60,2/200,2/320}{%
    \pgfmathsetmacro{\px}{1.0*\c*cos(\t)}\pgfmathsetmacro{\py}{0.5*\c*sin(\t)}%
    \pgfmathsetmacro{\gx}{-\px/1.0}\pgfmathsetmacro{\gy}{-4*\py/1.0}% -grad(x^2/1+4y^2)/... (dir. al centro)
    \pgfmathsetmacro{\gn}{sqrt(\gx*\gx+\gy*\gy)}%
    \pgfmathsetmacro{\ux}{0.62*\gx/\gn}\pgfmathsetmacro{\uy}{0.62*\gy/\gn}%
    \draw[vect, acento, line width=1pt] (\px,\py) -- ++(\ux,\uy);
    \fill (\px,\py) circle (1pt);}
  \node[font=\footnotesize, text=black!80, anchor=west] at (3.2,1.3) {curvas de nivel};
  
  \node[font=\footnotesize, text=azulprof, anchor=west] at (3.2,-1.3) {$h_3<h_2<h_1$};
\end{tikzpicture}\\[4pt]{\small\scshape (b)}\; $\nab h\perp\{h=\text{cte}\}$
\end{minipage}
\caption{Interpretación geométrica del gradiente. (a)~El incremento
$\dif f$ es el producto de $\norm{\nab f}$ por la proyección del desplazamiento
sobre la dirección del gradiente: es máximo si $\dl\parallel\nab f$ y nulo si
$\dl\perp\nab f$. (b)~Mapa topográfico de una colina: el gradiente de la altura
$h$ es perpendicular a las curvas de nivel y apunta hacia la cumbre (máxima
pendiente).}
\label{fig:gradiente-geo}
\end{figure}
